\chapter{轴对称}\label{ch:g6:symmetry}

沿一条直线折叠一张纸：每个点都会落在另一个点上，
即其镜像点。这种折叠称为\emph{轴对称}（或
反射），是本书学习的第一种几何变换——
第二种变换，中心对称，见\cref{ch:g7:central}。

\section{反射}

\begin{definition}[对称点]\label{def:g6:symmetry:def}
设 $d$ 为一条直线。所谓\emph{对称点}\index{对称点}，
是指：点 $M$ 关于 $d$ 的对称点为点 $M'$，使得 $d$ 为
$[MM']$ 的\omterm{def:g6:symmetry:bisector}{垂直平分线}：\omterm{def:g6:lines:objects}{线段} $[MM']$ 与
$d$ \omterm{def:g6:lines:perp}{垂直}，且 $d$ 经过它的中点。\emph{位于} $d$ 上的点
其对称点就是它自身。
\end{definition}

\begin{omfigure}
\begin{tikzpicture}[scale=1.05]
  \draw[thick, omExo] (0,-1.7) -- (0,1.9) node[above, font=\small] {$d$};
  \fill (-2.1,1.0) circle (1.6pt) node[left, font=\small] {$M$};
  \fill (2.1,1.0) circle (1.6pt) node[right, font=\small] {$M'$};
  \draw[dashed] (-2.1,1.0) -- (2.1,1.0);
  \draw (-0.10,0.90) -- (0.10,1.10);
  \draw (-1.15,0.90) -- (-0.95,1.10);
  \draw (0.95,0.90) -- (1.15,1.10);
  \draw[thin] (0,1.0) ++(0.16,0) -- ++(0,0.16) -- ++(-0.16,0);
  \fill (-1.3,-1.1) circle (1.6pt) node[left, font=\small] {$N$};
  \fill (1.3,-1.1) circle (1.6pt) node[right, font=\small] {$N'$};
  \draw[dashed] (-1.3,-1.1) -- (1.3,-1.1);
  \fill (0,0.2) circle (1.6pt) node[above right, font=\small] {$P = P'$};
\end{tikzpicture}

{\small 沿直线 $d$ 反射：\omterm{def:g6:symmetry:axis}{对称轴}与
$[MM']$ \omterm{def:g6:lines:perp}{垂直}，并经过其中点（等长记号）。点 $P$
位于\omterm{def:g6:symmetry:axis}{对称轴}上时不动。}
\end{omfigure}

\begin{method}[作对称点]\label{met:g6:symmetry:construct}
求 $M'$，即 $M$ 关于直线 $d$ 的\omterm{def:g6:symmetry:def}{对称点}：
\begin{enumerate}
  \item 作与 $d$ \omterm{def:g6:lines:perp}{垂直}且过 $M$ 的直线（用三角板）；记
        $H$ 为它与 $d$ 的交点；
  \item 用圆规或直尺量出 $MH$；
  \item 取点 $M'$，使其在同一垂线上且位于 $d$ 的另一侧，
        使 $HM' = MH$。
\end{enumerate}
在方格纸上作图更快：数出从 $M$ 到\omterm{def:g6:symmetry:axis}{对称轴}的方格数，
再在另一侧数出相同的格数。
\end{method}

\begin{example}[在方格纸上]\label{ex:g6:symmetry:grid}
若\omterm{def:g6:symmetry:axis}{对称轴}为竖直格线，且 $M$ 位于其左侧 $3$ 格，
则 $M'$ 在其右侧 $3$ 格处，与之同行。若\omterm{def:g6:symmetry:axis}{对称轴}为斜线，
则计数沿对角线进行——
始终与\omterm{def:g6:symmetry:axis}{对称轴}\omterm{def:g6:lines:perp}{垂直}。
\end{example}

\begin{omfigure}
\begin{tikzpicture}[scale=0.55]
  \draw[gray!40, thin] (0,0) grid (12,6);
  \draw[thick, omExo] (6,-0.4) -- (6,6.4);
  \draw[thick, omDef, fill=omDef!15]
    (1,1) -- (4,1) -- (4,3) -- (3,3) -- (3,5) -- (1,5) -- cycle;
  \draw[thick, omThm, fill=omThm!15]
    (11,1) -- (8,1) -- (8,3) -- (9,3) -- (9,5) -- (11,5) -- cycle;
  \node[omDef, font=\small] at (2.4,0.45) {图形};
  \node[omThm, font=\small] at (9.6,0.45) {其镜像};
\end{tikzpicture}

{\small 图形及其关于竖直\omterm{def:g6:symmetry:axis}{对称轴}的反射：每个
\omterm{def:g5:solids:def}{顶点}逐格反射，左右方向互换。}
\end{omfigure}

\begin{proposition}[反射保持的性质]\label{prop:g6:symmetry:preserved}
图形关于一条直线的对称图形，与原图形\emph{形状相同、
大小相同}：长度、角、\omterm{def:g4:measure:perimeter}{周长}与\omterm{def:g4:measure:area}{面积}都
保持不变。\omterm{def:g6:lines:objects}{线段}的对称图形是\omterm{def:g6:lines:objects}{线段}，圆的
对称图形是\omterm{def:g3:shapes:shapes}{半径}相同的圆，直线的对称图形是直线。
\end{proposition}
\admitted

\begin{remark}
有一项\emph{不}保持：方向。反射后的“b”
变成“d”——镜像书写。若描下你的图形，并翻转
描图纸，使其与像重合，则反射正确。
\end{remark}

\section{图形的对称轴}

\begin{definition}[对称轴]\label{def:g6:symmetry:axis}
直线 $d$ 称为图形的\emph{对称轴}\index{对称轴}，当关于
$d$ 的反射使该图形恰好与自身重合时，
——沿 $d$ 折叠，两半恰好重合。
\end{definition}

\begin{omfigure}
\begin{tikzpicture}[scale=0.62]
  % isosceles triangle
  \draw[thick] (0,0) -- (2.4,0) -- (1.2,2.6) -- cycle;
  \draw[omThm, dashed] (1.2,-0.5) -- (1.2,3.0);
  \node[below, font=\small] at (1.2,-0.7) {等腰：$1$ 条对称轴};
  % rectangle
  \begin{scope}[xshift=5.6cm]
  \draw[thick] (0,0.3) rectangle (3.2,2.3);
  \draw[omThm, dashed] (1.6,-0.4) -- (1.6,3.0);
  \draw[omThm, dashed] (-0.5,1.3) -- (3.7,1.3);
  \node[below, font=\small] at (1.6,-0.7) {长方形：$2$ 条对称轴};
  \end{scope}
  % square
  \begin{scope}[xshift=12.0cm]
  \draw[thick] (0.3,0.1) rectangle (2.7,2.5);
  \draw[omThm, dashed] (1.5,-0.5) -- (1.5,3.1);
  \draw[omThm, dashed] (-0.3,1.3) -- (3.3,1.3);
  \draw[omThm, dashed] (0.0,-0.2) -- (3.0,2.8);
  \draw[omThm, dashed] (0.0,2.8) -- (3.0,-0.2);
  \node[below, font=\small] at (1.5,-0.7) {正方形：$4$ 条对称轴};
  \end{scope}
\end{tikzpicture}

{\small 熟悉图形的\omterm{def:g6:symmetry:axis}{对称轴}。注意
长方形：其对角线\emph{不是}\omterm{def:g6:symmetry:axis}{对称轴}（沿对角线
折叠，两半不重合，除非它是正方形）。}
\end{omfigure}

\begin{example}\label{ex:g6:symmetry:shapes}
数一数\omterm{def:g6:symmetry:axis}{对称轴}的条数：
\begin{itemize}
  \item \omterm{def:g6:shapes:triangles}{等腰三角形}：$1$ 条（过\omterm{def:g5:solids:def}{顶点}与底边
        中点）；\omterm{def:g6:shapes:triangles}{等边三角形}：$3$ 条；
  \item 长方形：$2$ 条（两条“中线”，\emph{不是}
        对角线）；\omterm{def:g6:shapes:quadrilaterals}{菱形}：$2$ 条（它的对角线！）；正方形：$4$ 条；
  \item 圆：过圆心的每一条直线都是\omterm{def:g6:symmetry:axis}{对称轴}，有无穷多条。
\end{itemize}
\end{example}

\begin{example}[对称证明相等]\label{ex:g6:symmetry:proves}
为什么\omterm{def:g6:shapes:triangles}{等腰三角形}的两个底角相等？把三角形沿其\omterm{def:g6:symmetry:axis}{对称轴}
折叠：左半部分恰好与右半部分重合，
因此左底角恰好落在右底角上——
二者的度数必然相等
（由\cref{prop:g6:symmetry:preserved}可知，反射保持角的大小）。
\end{example}

\section{垂直平分线}

\begin{definition}[垂直平分线]\label{def:g6:symmetry:bisector}
称\omterm{def:g6:lines:objects}{线段} $[AB]$ 的\emph{垂直平分线}\index{垂直平分线}
为经过其中点且与 $[AB]$ \omterm{def:g6:lines:perp}{垂直}的直线，
它是交换 $A$ 与 $B$ 的\omterm{def:g6:symmetry:axis}{对称轴}。
\end{definition}

\begin{proposition}[等距性]\label{prop:g6:symmetry:bisector}
点 $M$ 位于 $[AB]$ 的\omterm{def:g6:symmetry:bisector}{垂直平分线}上，当且仅当它
到 $A$ 与 $B$ 的距离相等：$MA = MB$。
\end{proposition}
\admitted

\begin{method}[圆规作图]\label{met:g6:symmetry:bisector}
不借助测量，作\omterm{def:g6:lines:objects}{线段} $[AB]$ 的\omterm{def:g6:symmetry:bisector}{垂直平分线}：
\begin{enumerate}
  \item 圆规张开的宽度大于 $AB$ 的\omterm{def:g3:division:half}{一半}；
  \item 以 $A$ 为圆心，在\omterm{def:g6:lines:objects}{线段}上下两侧各画一段圆弧；
  \item 以 $B$ 为圆心，用\emph{相同}的张口画两段圆弧；
  \item 连接两个交点，这条直线即为
        \omterm{def:g6:symmetry:bisector}{垂直平分线}（每个交点到 $A$ 与 $B$ 的距离都相等）。
\end{enumerate}
\end{method}

\section{练习}

\begin{exercise}[$\star$]\label{exo:g6:symmetry:1}
在方格纸上画出：竖直直线 $d$，点 $A$ 位于
左侧两格处，点 $B$ 位于左侧五格、高一格处。
作出 $A'$ 与 $B'$，二者分别为它们关于 $d$ 的\omterm{def:g6:symmetry:def}{对称点}。
\end{exercise}

\begin{exercise}[$\star$]\label{exo:g6:symmetry:2}
画一条直线 $d$，并取一点 $M$，使其距 $3$ cm 于 $d$（不在
方格纸上）。用三角板和直尺作出 $M'$，即 $M$ 关于 $d$ 的\omterm{def:g6:symmetry:def}{对称点}，
并按照\cref{met:g6:symmetry:construct}作图。
$M$ 与 $M'$ 相距多远？
\end{exercise}

\begin{exercise}[$\star$]\label{exo:g6:symmetry:3}
在英文大写字母中，取最简形式书写时，哪些字母
有竖直\omterm{def:g6:symmetry:axis}{对称轴}（如 A）？有水平\omterm{def:g6:symmetry:axis}{对称轴}（如 B）？
两种都有的呢？
\end{exercise}

\begin{exercise}[$\star$]\label{exo:g6:symmetry:4}
下列图形各有几条\omterm{def:g6:symmetry:axis}{对称轴}：\omterm{def:g6:shapes:triangles}{等边三角形}？长方形
（不是正方形）？\omterm{def:g6:shapes:quadrilaterals}{菱形}（不是正方形）？正方形？圆？
请画出每个图形及其\omterm{def:g6:symmetry:axis}{对称轴}。
\end{exercise}

\begin{exercise}[$\star$]\label{exo:g6:symmetry:5}
判断下列说法是否正确，并说明理由：“长方形的对角线是该长方形的
\omterm{def:g6:symmetry:axis}{对称轴}。”可用折叠论证或作图说明。
\end{exercise}

\begin{exercise}[$\star$]\label{exo:g6:symmetry:6}
三角形 $ABC$ 关于直线 $d$ 作反射得到 $A'B'C'$。已知
$AB = 4$ cm，$\widehat{BAC} = 70^\circ$，三角形 $ABC$ 的\omterm{def:g4:measure:perimeter}{周长}为
$12$ cm。不作图，直接写出 $A'B'$、
$\widehat{B'A'C'}$ 及 $A'B'C'$ 的\omterm{def:g4:measure:perimeter}{周长}。
\end{exercise}

\begin{exercise}[$\star$]\label{exo:g6:symmetry:7}
画一条\omterm{def:g6:lines:objects}{线段} $[AB]$，长 $7$ cm，并用圆规作出它的
\omterm{def:g6:symmetry:bisector}{垂直平分线}（见\cref{met:g6:symmetry:bisector}）。在其上任取
一点 $M$，用直尺检验 $MA = MB$。
\end{exercise}

\begin{exercise}[$\star\star$]\label{exo:g6:symmetry:8}
平地花园中的两棵树分别位于点 $A$ 与 $B$ 处，
在何处设置一个喷泉，才能使其到两棵树的距离相等？
请描述\emph{全部}可能的位置。
\end{exercise}

\begin{exercise}[$\star\star$]\label{exo:g6:symmetry:9}
在方格纸上画出\omterm{def:g6:symmetry:axis}{对称轴} $d$（方格的一条对角线，
与水平线成 $45^\circ$），在其一侧画一个小的 L 形图形，作出它
关于 $d$ 的对称图形。（逐个\omterm{def:g5:solids:def}{顶点}作反射：对 $45^\circ$
轴，轴右侧 $2$ 格的点，对称后落在它的上方
$2$ 格处。）
\end{exercise}

\begin{exercise}[$\star\star$]\label{exo:g6:symmetry:10}
画一个圆心为 $O$ 的圆，再画一条直线 $d$ 经过 $O$，并在圆上取一点 $M$。
$M$ 关于 $d$ 的\omterm{def:g6:symmetry:def}{对称点}在何处？说明
为什么反射把该圆映到它自身。
\end{exercise}

\begin{exercise}[$\star\star$]\label{exo:g6:symmetry:11}
利用等距性质（\cref{prop:g6:symmetry:bisector}），
说明为什么
在\cref{met:g6:symmetry:bisector}中的两段圆弧的两个交点
确实位于 $[AB]$ 的\omterm{def:g6:symmetry:bisector}{垂直平分线}上。
\end{exercise}

\begin{exercise}[$\star\star\star$]\label{exo:g6:symmetry:12}
一个台球位于点 $A$ 处，须经直墙 $d$ 反弹后
到达点 $B$（两点位于墙的同一侧）。作 $B$
关于 $d$ 的\omterm{def:g6:symmetry:def}{对称点} $B'$，并画出\omterm{def:g6:lines:objects}{线段} $[AB']$。说明为什么
最佳反弹点正是 $[AB']$ 与墙的交点。（思路：路径
$A \to P \to B$ 的长度与 $A \to P \to B'$ 相等。）
\end{exercise}

\section{问题：喷泉、破碎的盘子
与走动的镜子}

\begin{problem}[{周末问题——三角形的三条垂直平分线
交于一点：过三点的圆，
以及两面镜子共同完成的事}]
\label{pb:g6:symmetry:1}
在\cref{exo:g6:symmetry:8}中，喷泉须与
\emph{两}棵树距离相等，答案是一整条
线上的点。本题又种下\emph{第三}棵树，于是整条
线便收缩为唯一确定的一点。这
一点蕴含着几何的宝藏：它能让你画出唯一的
一个过给定三点的圆，还能由一块碎片
复原整个盘子，而且它属于每一个三角形。最后
一部分用两面镜子做了一个实验，其结果
预示了下一学年的几何内容。

\medskip
\noindent\textbf{第一部分——两棵树，然后是一张
领地分布图。}
\begin{enumerate}
  \item 画两点 $A$ 和 $B$，使 $AB = 6$ cm，并
        用圆规作 $[AB]$ 的\omterm{def:g6:symmetry:bisector}{垂直平分线}
        （见\cref{met:g6:symmetry:bisector}）。在其上取一点
        $M$，用直尺检验 $MA = MB$。
  \item \Cref{prop:g6:symmetry:bisector}同时说明两件事：
        一是\omterm{def:g6:symmetry:bisector}{垂直平分线}\emph{上}的每一点
        到 $A$、$B$ 等距；二是每个等距点都\emph{在}
        \omterm{def:g6:symmetry:bisector}{垂直平分线}上。两者之中哪一个回答了
        在\cref{exo:g6:symmetry:8}中提出的喷泉问题？
        命题对\emph{不}在
        \omterm{def:g6:symmetry:bisector}{垂直平分线}上的点说了什么？
  \item 取一点 $M$，使它与
        $A$ 位于\omterm{def:g6:symmetry:bisector}{垂直平分线}的同一侧。\omterm{def:g6:lines:objects}{线段} $[MB]$ 与\omterm{def:g6:symmetry:bisector}{垂直平分线}交于点
        $P$。比较路径 $M \to P \to B$ 与
        路径 $M \to P \to A$，并说明为什么 $M$ 离 $A$ 较近，
        而非离 $B$ 更近：\omterm{def:g6:symmetry:bisector}{垂直平分线}把纸面分成“
        $A$ 侧”与“$B$ 侧”两部分。
  \item 两所学校 $A$ 与 $B$ 服务同一座镇。每个孩子都走
        到离自己较近的学校。画一张小地图，标出 $A$ 与 $B$，
        并涂出镇上由学校
        $A$ 服务的确切区域。哪条线构成边界？
  \item 再开设第三所学校 $C$（选址使三所学校
        构成一个三角形）。分别作
        $[AB]$ 与 $[BC]$ 的\omterm{def:g6:symmetry:bisector}{垂直平分线}，并
        记它们的交点为 $O$。$O$ 满足哪两个
        距离相等关系，为什么？
\end{enumerate}

\medskip
\noindent\textbf{第二部分——三棵树、一点与一个
圆。}
\begin{enumerate}[resume]
  \item 由第 5 题可推得 $OA = OC$ —— 因而
        $O$ 也位于 $[AC]$ 的\omterm{def:g6:symmetry:bisector}{垂直平分线}上，
        而你从未画过它。由此得出：\emph{三角形各边的三条
        \omterm{def:g6:symmetry:bisector}{垂直平分线}
        都经过同一点。}
  \item 说明为什么以 $O$ 为圆心、$OA$ 为\omterm{def:g3:shapes:shapes}{半径}的圆
        同样经过 $B$ 与 $C$：它是
        \emph{这个}过三点的圆——
        \emph{外接圆}\index{外接圆}，
        即三角形 $ABC$ 的\omterm{pb:g6:symmetry:1}{外接圆}。
  \item 作一个三角形，三边分别为 $6$ cm、$5$ cm 和
        $4$ cm（见\cref{met:g6:shapes:construct}），再作它的
        \omterm{pb:g6:symmetry:1}{外接圆}。为什么只需画
        \emph{两}条\omterm{def:g6:symmetry:bisector}{垂直平分线}就能找到圆心？
  \item 破碎的盘子：考古学家挖出一块碎片，其
        唯一完好的部分是盘子圆形边缘的一段。
        请描述一种方法，恢复原盘子的圆心与
        \omterm{def:g3:shapes:shapes}{半径}，并加以检验：用圆规画一段圆弧，
        “忘掉”圆心，
        在圆弧上取三点，并重建这个圆。
  \item 三点是否可能\emph{没有}过它们的圆？
        假设 $A$、$B$、$C$ 共线。
        $[AB]$ 的\omterm{def:g6:symmetry:bisector}{垂直平分线}与
        $[BC]$ 的\omterm{def:g6:symmetry:bisector}{垂直平分线}能说明什么（见\cref{prop:g6:lines:facts}）？由此得出结论。
\end{enumerate}

\medskip
\noindent\textbf{第三部分——处处是\omterm{def:g6:symmetry:axis}{对称轴}，以及走动的
镜子。}
\begin{enumerate}[resume]
  \item 在\cref{ex:g6:symmetry:shapes}中数出：\omterm{def:g6:shapes:triangles}{等边三角形}有 $3$ 条
        \omterm{def:g6:symmetry:axis}{对称轴}，正方形有 $4$ 条。猜想
        正\omterm{def:g4:geometry:polygon}{五边形}与正\omterm{def:g4:geometry:polygon}{六边形}的\omterm{def:g6:symmetry:axis}{对称轴}条数，
        并描述各自\omterm{def:g6:symmetry:axis}{对称轴}经过的位置（边数为\omterm{def:g3:numbers:evenodd}{奇数}
        与\omterm{def:g3:numbers:evenodd}{偶数}相比，两种情形
        有所不同）。
  \item 将一张纸对折一次，沿两层剪出一个图形；
        展开。为什么孔洞\emph{总}有\omterm{def:g6:symmetry:axis}{对称轴}，
        无论剪的是什么形状？
        这条\omterm{def:g6:symmetry:axis}{对称轴}是哪一条直线？
  \item 将纸对折两次，第二次折痕与第一次
        \omterm{def:g6:lines:perp}{垂直}，剪出一个图形后展开。孔洞至少有几条
        \omterm{def:g6:symmetry:axis}{对称轴}，以及能看到多少份剪出的图形？
        （纸雪花正是
        这一想法的进一步折叠。）
  \item 走动的镜子：在方格纸上画一条竖直的\omterm{def:g6:symmetry:axis}{对称轴}
        $d_1$，在其左侧画一面小旗 $F$，并画第二条
        竖直\omterm{def:g6:symmetry:axis}{对称轴} $d_2$，距 $4$ 格，在 $d_1$ 的右侧。
        作 $F$ 关于 $d_1$ 的对称图形 $F'$，再
        作 $F'$ 关于 $d_2$ 的对称图形 $F''$。比较 $F$ 与 $F''$：
        方向相同，还是互为镜像？图形移动了多少
        格？与两轴间的距离相比，
        有何关系？
  \item 做同样的实验，但令 $d_2$ \emph{\omterm{def:g6:lines:perp}{垂直}}于
        $d_1$：先将小旗关于竖直轴作对称，
        再将所得像关于水平轴作对称。描述 $F''$
        相对于 $F$ 的位置，以及它与两轴的
        交点的位置。（你刚刚发现了下一学年的
        变换，见\cref{ch:g7:central}。）
\end{enumerate}
\end{problem}
